\documentclass[12pt,letterpaper]{article}
\usepackage{graphicx,textcomp}
\usepackage{natbib}
\usepackage{setspace}
\usepackage{fullpage}
\usepackage{color}
\usepackage[reqno]{amsmath}
\usepackage{amsthm}
\usepackage{fancyvrb}
\usepackage{amssymb,enumerate}
\usepackage[all]{xy}
\usepackage{endnotes}
\usepackage{lscape}
\newtheorem{com}{Comment}
\usepackage{float}
\usepackage{hyperref}
\newtheorem{lem} {Lemma}
\newtheorem{prop}{Proposition}
\newtheorem{thm}{Theorem}
\newtheorem{defn}{Definition}
\newtheorem{cor}{Corollary}
\newtheorem{obs}{Observation}
\usepackage[compact]{titlesec}
\usepackage{dcolumn}
\usepackage{tikz}
\usetikzlibrary{arrows}
\usepackage{multirow}
\usepackage{xcolor}
\newcolumntype{.}{D{.}{.}{-1}}
\newcolumntype{d}[1]{D{.}{.}{#1}}
\definecolor{light-gray}{gray}{0.65}
\usepackage{url}
\usepackage{listings}
\usepackage{color}

\definecolor{codegreen}{rgb}{0,0.6,0}
\definecolor{codegray}{rgb}{0.5,0.5,0.5}
\definecolor{codepurple}{rgb}{0.58,0,0.82}
\definecolor{backcolour}{rgb}{0.95,0.95,0.92}

\lstdefinestyle{mystyle}{
	backgroundcolor=\color{backcolour},   
	commentstyle=\color{codegreen},
	keywordstyle=\color{magenta},
	numberstyle=\tiny\color{codegray},
	stringstyle=\color{codepurple},
	basicstyle=\footnotesize,
	breakatwhitespace=false,         
	breaklines=true,                 
	captionpos=b,                    
	keepspaces=true,                 
	numbers=left,                    
	numbersep=5pt,                  
	showspaces=false,                
	showstringspaces=false,
	showtabs=false,                  
	tabsize=2
}
\lstset{style=mystyle}
\newcommand{\Sref}[1]{Section~\ref{#1}}
\newtheorem{hyp}{Hypothesis}

\title{Problem Set 3}
\date{Due: March 25, 2026}
\author{Applied Stats II}


\begin{document}
	\maketitle
	\section*{Instructions}
	\begin{itemize}
	\item Please show your work! You may lose points by simply writing in the answer. If the problem requires you to execute commands in \texttt{R}, please include the code you used to get your answers. Please also include the \texttt{.R} file that contains your code. If you are not sure if work needs to be shown for a particular problem, please ask.
\item Your homework should be submitted electronically on GitHub in \texttt{.pdf} form.
\item This problem set is due before 23:59 on Wednesday March 25, 2026. No late assignments will be accepted.
	\end{itemize}

	\vspace{.25cm}
\section*{Question 1}
\vspace{.25cm}
\noindent We are interested in how governments' management of public resources impacts economic prosperity. Our data come from \href{https://www.researchgate.net/profile/Adam_Przeworski/publication/240357392_Classifying_Political_Regimes/links/0deec532194849aefa000000/Classifying-Political-Regimes.pdf}{Alvarez, Cheibub, Limongi, and Przeworski (1996)} and is labelled \texttt{gdpChange.csv} on GitHub. The dataset covers 135 countries observed between 1950 or the year of independence or the first year forwhich data on economic growth are available ("entry year"), and 1990 or the last year for which data on economic growth are available ("exit year"). The unit of analysis is a particular country during a particular year, for a total $>$ 3,500 observations. 

\begin{itemize}
	\item
	Response variable: 
	\begin{itemize}
		\item \texttt{GDPWdiff}: Difference in GDP between year $t$ and $t-1$. Possible categories include: "positive", "negative", or "no change"
	\end{itemize}
	\item
	Explanatory variables: 
	\begin{itemize}
		\item
		\texttt{REG}: 1=Democracy; 0=Non-Democracy
		\item
		\texttt{OIL}: 1=if the average ratio of fuel exports to total exports in 1984-86 exceeded 50\%; 0= otherwise
	\end{itemize}
	
\end{itemize}
\newpage
\noindent Please answer the following questions:

\begin{enumerate}
	\item Construct and interpret an unordered multinomial logit with \texttt{GDPWdiff} as the output and "no change" as the reference category, including the estimated cutoff points and coefficients.
	\item Construct and interpret an ordered multinomial logit with \texttt{GDPWdiff} as the outcome variable, including the estimated cutoff points and coefficients.
	
	
\end{enumerate}

\section*{Question 2} 
\vspace{.25cm}

\noindent Consider the data set \texttt{MexicoMuniData.csv}, which includes municipal-level information from Mexico. The outcome of interest is the number of times the winning PAN presidential candidate in 2006 (\texttt{PAN.visits.06}) visited a district leading up to the 2009 federal elections, which is a count. Our main predictor of interest is whether the district was highly contested, or whether it was not (the PAN or their opponents have electoral security) in the previous federal elections during 2000 (\texttt{competitive.district}), which is binary (1=close/swing district, 0="safe seat"). We also include \texttt{marginality.06} (a measure of poverty) and \texttt{PAN.governor.06} (a dummy for whether the state has a PAN-affiliated governor) as additional control variables. 

\begin{enumerate}
	\item [(a)]
	Run a Poisson regression because the outcome is a count variable. Is there evidence that PAN presidential candidates visit swing districts more? Provide a test statistic and p-value.

	\item [(b)]
	Interpret the \texttt{marginality.06} and \texttt{PAN.governor.06} coefficients.
	
	\item [(c)]
	Provide the estimated mean number of visits from the winning PAN presidential candidate for a hypothetical district that was competitive (\texttt{competitive.district}=1), had an average poverty level (\texttt{marginality.06} = 0), and a PAN governor (\texttt{PAN.governor.06}=1).
	
\end{enumerate}
\newpage
\section*{My Answers}
\subsection*{Question 1}

\noindent \textbf{(1) Unordered multinomial logit.}
To begin, I estimated an unordered multinomial logit model in which \texttt{GDPWdiff} was recoded into three categories: \texttt{negative}, \texttt{no change}, and \texttt{positive}. Following the instructions, \texttt{no change} was used as the reference category. The model therefore estimates two comparisons: \texttt{negative} versus \texttt{no change}, and \texttt{positive} versus \texttt{no change}, as a function of \texttt{REG} and \texttt{OIL}.\\

% Table created by stargazer v.5.2.3 by Marek Hlavac, Social Policy Institute. E-mail: marek.hlavac at gmail.com
% Date and time: Thu, Mar 12, 2026 - 12:21:02
\begin{table}[!htbp] \centering 
	\caption{Unordered Multinomial Logit Model} 
	\label{} 
	\begin{tabular}{@{\extracolsep{5pt}}lcc} 
		\\[-1.8ex]\hline 
		\hline \\[-1.8ex] 
		& \multicolumn{2}{c}{\textit{Dependent variable:}} \\ 
		\cline{2-3} 
		\\[-1.8ex] & GDPWdiff (ref. = no change) & positive \\ 
		& negative vs no change & positive vs no change \\ 
		\\[-1.8ex] & (1) & (2)\\ 
		\hline \\[-1.8ex] 
		REG & 1.379$^{*}$ & 1.769$^{**}$ \\ 
		& (0.769) & (0.767) \\ 
		& & \\ 
		OIL & 4.784 & 4.576 \\ 
		& (6.885) & (6.885) \\ 
		& & \\ 
		Constant & 3.805$^{***}$ & 4.534$^{***}$ \\ 
		& (0.271) & (0.269) \\ 
		& & \\ 
		\hline \\[-1.8ex] 
		Akaike Inf. Crit. & 4,690.770 & 4,690.770 \\ 
		\hline 
		\hline \\[-1.8ex] 
		\textit{Note:}  & \multicolumn{2}{r}{$^{*}$p$<$0.1; $^{**}$p$<$0.05; $^{***}$p$<$0.01} \\ 
	\end{tabular} 
\end{table} 

\noindent The coefficient on \texttt{REG} is positive in both comparisons. For \texttt{negative} versus \texttt{no change}, the estimate is positive but only marginally statistically significant. For \texttt{positive} versus \texttt{no change}, the coefficient is positive and statistically significant at the 5\% level. Substantively, this suggests that when a country is democratic rather than non-democratic, it has higher log-odds of falling into either the negative or positive GDP-change category rather than the \texttt{no change} category, holding \texttt{OIL} constant. However, the positive comparison is estimated more precisely than the negative one.\\

\noindent The coefficient on \texttt{OIL} is positive in both logits, but in neither case is it statistically significant. Moreover, the standard errors are very large, indicating substantial uncertainty around these estimates. A key reason for this instability is that the reference category, \texttt{no change}, contains only 16 observations. Since both multinomial comparisons are made against this very small baseline group, the resulting coefficients should be interpreted with caution.\\

\noindent \textbf{(2) Ordered multinomial logit.}
I next estimated an ordered multinomial logit model, treating the outcome as ordinal with the ranking \texttt{negative < no change < positive}. Unlike the unordered model, this specification uses the natural ordering of the outcome and estimates how the predictors affect the odds of being in a higher rather than lower GDP-change category.\\

% Table created by stargazer v.5.2.3 by Marek Hlavac, Social Policy Institute. E-mail: marek.hlavac at gmail.com
% Date and time: Thu, Mar 12, 2026 - 12:21:02
\begin{table}[!htbp] \centering 
	\caption{Ordered Multinomial Logit Model} 
	\label{} 
	\begin{tabular}{@{\extracolsep{5pt}}lc} 
		\\[-1.8ex]\hline 
		\hline \\[-1.8ex] 
		& \multicolumn{1}{c}{\textit{Dependent variable:}} \\ 
		\cline{2-2} 
		\\[-1.8ex] & GDPWdiff (negative < no change < positive) \\ 
		\hline \\[-1.8ex] 
		REG & 0.398$^{***}$ \\ 
		& (0.075) \\ 
		& \\ 
		OIL & $-$0.199$^{*}$ \\ 
		& (0.116) \\ 
		& \\ 
		\hline \\[-1.8ex] 
		Observations & 3,721 \\ 
		\hline 
		\hline \\[-1.8ex] 
		\textit{Note:}  & \multicolumn{1}{r}{$^{*}$p$<$0.1; $^{**}$p$<$0.05; $^{***}$p$<$0.01} \\ 
	\end{tabular} 
\end{table} 

\noindent The coefficient on \texttt{REG} is positive and highly statistically significant. This indicates that, holding \texttt{OIL} constant, democratic countries have higher odds of being in a more favourable GDP-change category than non-democratic countries. The estimated odds ratio is approximately 1.49, which implies that democracy is associated with about a 49\% increase in the cumulative odds of being in a higher outcome category.\\

\noindent The coefficient on \texttt{OIL} is negative, implying that higher oil dependence is associated with lower odds of being in a better GDP-change category. The estimated odds ratio is approximately 0.82, corresponding to a reduction of about 18\% in the cumulative odds of being in a higher category. However, this effect is only marginally significant at the 10\% level and is not statistically significant at the conventional 5\% threshold.\\

\noindent The ordered logit also estimates two cut points, separating \texttt{negative} from \texttt{no change} and \texttt{no change} from \texttt{positive}. These thresholds are very close together, which likely reflects the fact that the middle category is extremely small.\\

\noindent \textbf{Conclusion.}
Taken together, the ordered multinomial logit provides the clearer substantive result. It suggests that democratic countries are more likely to experience more favourable GDP-change outcomes, whereas the evidence for an effect of oil dependence is weaker. By contrast, the unordered multinomial logit is more difficult to interpret confidently because both comparisons rely on the very small \texttt{no change} reference category.\\

\noindent\textbf{R code used:}

\lstinputlisting[language=R, firstline=35, lastline=124]{PS3_jl.R}
\newpage

\subsection*{Question 2}

\textbf{Poisson regression model.}
Because \texttt{PAN.visits.06} is a count variable, I estimated a Poisson regression model with \texttt{competitive.district}, \texttt{marginality.06}, and \texttt{PAN.governor.06} as predictors.

% Table created by stargazer v.5.2.3 by Marek Hlavac, Social Policy Institute. E-mail: marek.hlavac at gmail.com
% Date and time: Thu, Mar 12, 2026 - 13:10:01
\begin{table}[!htbp] \centering 
	\caption{Poisson Regression Model} 
	\label{} 
	\begin{tabular}{@{\extracolsep{5pt}}lc} 
		\\[-1.8ex]\hline 
		\hline \\[-1.8ex] 
		& \multicolumn{1}{c}{\textit{Dependent variable:}} \\ 
		\cline{2-2} 
		\\[-1.8ex] & PAN candidate visits in 2006 \\ 
		\hline \\[-1.8ex] 
		competitive.district & $-$0.081 \\ 
		& (0.171) \\ 
		& \\ 
		marginality.06 & $-$2.080$^{***}$ \\ 
		& (0.117) \\ 
		& \\ 
		PAN.governor.06 & $-$0.312$^{*}$ \\ 
		& (0.167) \\ 
		& \\ 
		Constant & $-$3.810$^{***}$ \\ 
		& (0.222) \\ 
		& \\ 
		\hline \\[-1.8ex] 
		Observations & 2,407 \\ 
		Log Likelihood & $-$645.606 \\ 
		Akaike Inf. Crit. & 1,299.213 \\ 
		\hline 
		\hline \\[-1.8ex] 
		\textit{Note:}  & \multicolumn{1}{r}{$^{*}$p$<$0.1; $^{**}$p$<$0.05; $^{***}$p$<$0.01} \\ 
	\end{tabular} 
\end{table} 

\noindent \textbf{(a) Competitive districts.}
To assess whether PAN candidates visited competitive districts more often in 2006, I examined the coefficient on \texttt{competitive.district}. The estimated coefficient is \(-0.081\), with a \(z\)-statistic of \(-0.477\) and a \(p\)-value of \(0.634\). This indicates that, holding the other covariates constant, there is no statistically significant evidence that competitive districts received more PAN candidate visits. The estimated effect is slightly negative rather than positive. The corresponding incidence-rate ratio is \(e^{-0.081} \approx 0.92\), implying that competitive districts had an expected visit count about 8\% lower than non-competitive districts. However, this difference is not statistically distinguishable from zero.\\

\noindent \textbf{(b) Interpretation of \texttt{marginality.06} and \texttt{PAN.governor.06}.}
The coefficient on \texttt{marginality.06} is \(-2.080\) and is highly statistically significant (\(p < 0.001\)). This implies that higher values of \texttt{marginality.06} are associated with fewer PAN candidate visits, holding the other predictors constant. Exponentiating the coefficient gives an incidence-rate ratio of \(e^{-2.080} \approx 0.125\). Thus, a one-unit increase in \texttt{marginality.06} is associated with the expected number of visits being multiplied by about 0.125, corresponding to an approximate 87.5\% decrease in the expected count.

\noindent The coefficient on \texttt{PAN.governor.06} is \(-0.312\), with a \(p\)-value of \(0.062\). This suggests that municipalities in states with a PAN governor tended to receive fewer PAN candidate visits than those in states without a PAN governor, holding the other variables constant. The corresponding incidence-rate ratio is \(e^{-0.312} \approx 0.73\), implying roughly a 27\% lower expected number of visits. However, this coefficient is only marginally significant at the 10\% level and does not reach the conventional 5\% threshold.\\

\noindent \textbf{(c) Predicted number of visits.}
For a district with \texttt{competitive.district = 1}, \texttt{marginality.06 = 0}, and \texttt{PAN.governor.06 = 1}, the model predicts an expected mean of \(0.01495\) PAN candidate visits. Rounded, this is approximately \(0.015\) visits. This is an estimated expected count rather than a predicted observed number of visits.\\

\noindent \textbf{Conclusion.}
Overall, the Poisson model provides no evidence that PAN candidates visited competitive districts more often in 2006. Instead, the clearest result is that higher marginality is associated with substantially fewer visits, while the negative association for municipalities in states with a PAN governor is weaker and only marginally significant.\\

\noindent\textbf{R code used:}

\lstinputlisting[language=R, firstline=147, lastline=217]{PS3_jl.R}

\end{document}
